\clearpage
\section{A 4 potential}
\subsection{finding the fields}\label{q2s1}
We'll find the fields tensor instead and derive the fields from it. as we know, the fields tensor is:
\begin{equation}
  \tensor{F}{^\mu^\nu} = \partial^\mu A^\nu-\partial^\nu A^\mu
\end{equation}
so
\begin{align}
  E_i &= \tensor{F}{_0_i}=\partial_0A_i-\partial_iA_0\\
  &= \partial_0\tensor{\eta}{_i_\alpha}A^\alpha-\partial_i\tensor{\eta}{_0_\beta}A^\beta\\
  &= -\partial_0A^i-\partial_iA^0\\
  &= -\diffp*{\ins{x^i e^{-\alpha((x^0)^2+(x^1)^2+(x^2)^2+(x^3)^2)}}}{x^0}-\diffp*{\ins{x^0 e^{-\alpha((x^0)^2+(x^1)^2+(x^2)^2+(x^3)^2)}}}{x^i}\\
  &= 2\alpha x^0x^ie^{-\alpha((x^0)^2+(x^1)^2+(x^2)^2+(x^3)^2)}+2\alpha x^0x^ie^{-\alpha((x^0)^2+(x^1)^2+(x^2)^2+(x^3)^2)}\\
  &= 4\alpha x^0x^ie^{-\alpha((x^0)^2+(x^1)^2+(x^2)^2+(x^3)^2)}
\end{align}
and
\begin{align}
	B_1 &= \tensor{F}{_3_2} = \partial_3A_2-\partial_2A_3\\
	&= \partial_3\tensor{\eta}{_2_2}A^2-\partial_2\tensor{\eta}{_3_3}A^3\\
	&= \partial_2A^3-\partial_3A^2\\
	&= \partial_2\ins{x^3e^{-\alpha((x^0)^2+(x^1)^2+(x^2)^2+(x^3)^2)}}-\partial_3\ins{x^2e^{-\alpha((x^0)^2+(x^1)^2+(x^2)^2+(x^3)^2)}}\\
	&= -2\alpha x^3x^2e^{-\alpha((x^0)^2+(x^1)^2+(x^2)^2+(x^3)^2)}+2\alpha x^2x^3e^{-\alpha((x^0)^2+(x^1)^2+(x^2)^2+(x^3)^2)}\\
	&= 0
\end{align}
from symmetry to choice of index,
\begin{align}
	B_2=B_3=0
\end{align}
so
\begin{align}
	\bold E=4\alpha x^0e^{-\alpha((x^0)^2+(x^1)^2+(x^2)^2+(x^3)^2)}\bold x && \bold B=0
\end{align}
\subsection{4-current}
Ee saw a relation between the 4-current and the fields tensor
\begin{equation}
  \frac{4\pi}{c}\bold J^\nu = \partial_\mu\tensor{F}{^\mu^\nu}
\end{equation}
so
\begin{align}
	J^0&=\frac{c}{4\pi}\partial_\mu\tensor{F}{^\mu^0}\\
	&= \frac{c}{4\pi}\partial_i\tensor{F}{^i^0}\\
	&= \frac{c}{4\pi}\partial_i\ins{4\alpha x^0x^ie^{-\alpha((x^0)^2+(x^1)^2+(x^2)^2+(x^3)^2)}}\\
	&= \frac{c}{\pi}\alpha x^0\ins{3e^{-\alpha((x^0)^2+(x^1)^2+(x^2)^2+(x^3)^2)}-2\alpha x^ix_ie^{-\alpha((x^0)^2+(x^1)^2+(x^2)^2+(x^3)^2)}}\\
	&= \frac{\alpha c}{\pi}x^0\ins{3-2\alpha x^ix_i}e^{-\alpha((x^0)^2+(x^1)^2+(x^2)^2+(x^3)^2)}
\end{align}
and
\begin{align}
	J^i&= \frac{c}{4\pi}\partial_\mu\tensor{F}{^\mu^i}\\
	&= \frac{c}{4\pi}\partial_0\tensor{F}{^0^i}\\
	&= -\frac{c}{4\pi}\partial_0\ins{4\alpha x^0x^ie^{-\alpha((x^0)^2+(x^1)^2+(x^2)^2+(x^3)^2)}}\\
	&= -\frac{\alpha c}{\pi}\ins{x^ie^{-\alpha((x^0)^2+(x^1)^2+(x^2)^2+(x^3)^2)}-2\alpha x^0x_0x^ie^{-\alpha((x^0)^2+(x^1)^2+(x^2)^2+(x^3)^2)}}\\
	&= -\frac{\alpha c}{\pi}x^i\ins{1-2\alpha x^0x_0}e^{-\alpha((x^0)^2+(x^1)^2+(x^2)^2+(x^3)^2)}
\end{align}
therefore
\begin{equation}
  \bold J^\mu = \frac{\alpha c}{\pi}\mqty(3-2\alpha x^ix_i\\2\alpha x^0x_0-1\\2\alpha x^0x_0-1\\2\alpha x^0x_0-1)\bold xe^{-\alpha((x^0)^2+(x^1)^2+(x^2)^2+(x^3)^2)}
\end{equation}
We will show conservation of charge exists by showing the continuity equation works. explicitly the 4-current is
\begin{equation}
  \bold J^\mu = \frac{\alpha c}{\pi}\mqty(3x^0-2\alpha x^0x^ix_i\\2\alpha x^0x_0x^1-x^1\\2\alpha x^0x_0x^2-x^2\\2\alpha x^0x_0x^3-x^3)\bold x e^{-\alpha((x^0)^2+(x^1)^2+(x^2)^2+(x^3)^2)}
\end{equation}
so
\begin{align}
	\partial_\mu\bold J^\mu &= \frac{\alpha c}{\pi}\diffp*{}{x^\mu}\mqty(3x^0-2\alpha x^0x^ix_i\\2\alpha x^0x_0x^1-x^1\\2\alpha x^0x_0x^2-x^2\\2\alpha x^0x_0x^3-x^3)\bold x e^{-\alpha((x^0)^2+(x^1)^2+(x^2)^2+(x^3)^2)}\\
	&= \frac{\alpha c}{\pi}\mqty(\partial_0&\partial_1&\partial_2&\partial_3)\mqty(3x^0-2\alpha x^0x^ix_i\\2\alpha x^0x_0x^1-x^1\\2\alpha x^0x_0x^2-x^2\\2\alpha x^0x_0x^3-x^3)\bold x e^{-\alpha((x^0)^2+(x^1)^2+(x^2)^2+(x^3)^2)}\\
	&= \frac{\alpha c}{\pi}\ins{3-2\alpha x^ix_i}\ins{1-2\alpha \ins{x^0}^2}e^{-\alpha((x^0)^2+(x^1)^2+(x^2)^2+(x^3)^2)}\\
	&+\frac{\alpha c}{\pi}\ins{2\alpha x^0x_0-1}\ins{1-2\alpha \ins{x^1}^2}e^{-\alpha((x^0)^2+(x^1)^2+(x^2)^2+(x^3)^2)}\\
	&+\frac{\alpha c}{\pi}\ins{2\alpha x^0x_0-1}\ins{1-2\alpha \ins{x^2}^2}e^{-\alpha((x^0)^2+(x^1)^2+(x^2)^2+(x^3)^2)}\\
	&+\frac{\alpha c}{\pi}\ins{2\alpha x^0x_0-1}\ins{1-2\alpha \ins{x^3}^2}e^{-\alpha((x^0)^2+(x^1)^2+(x^2)^2+(x^3)^2)}\\
  &= \frac{\alpha c}{\pi}\ins{2\alpha x^ix_i}\ins{1-2\alpha \ins{x^0}^2}e^{-\alpha((x^0)^2+(x^1)^2+(x^2)^2+(x^3)^2)}\\
	&+\frac{\alpha c}{\pi}\ins{2\alpha x^0x_0-1}\ins{-2\alpha \ins{x^1}^2}e^{-\alpha((x^0)^2+(x^1)^2+(x^2)^2+(x^3)^2)}\\
	&+\frac{\alpha c}{\pi}\ins{2\alpha x^0x_0-1}\ins{-2\alpha \ins{x^2}^2}e^{-\alpha((x^0)^2+(x^1)^2+(x^2)^2+(x^3)^2)}\\
	&+\frac{\alpha c}{\pi}\ins{2\alpha x^0x_0-1}\ins{-2\alpha \ins{x^3}^2}e^{-\alpha((x^0)^2+(x^1)^2+(x^2)^2+(x^3)^2)}\\
	&= 0
\end{align}
\subsection{Temporal Gauge}
We want to find a gauge transformation $\Lambda$ such that $A^0+\partial^0\Lambda=0$. Explicitly,
\begin{align}
	x^0e^{-\alpha((x^0)^2+(x^1)^2+(x^2)^2+(x^3)^2)}+\diffp{\Lambda}{x^0} = 0
\end{align}
therefore
\begin{equation}
  \Lambda = -\int \dd x^0 x^0e^{-\alpha((x^0)^2+(x^1)^2+(x^2)^2+(x^3)^2)}=\frac{1}{2\alpha}e^{-\alpha((x^0)^2+(x^1)^2+(x^2)^2+(x^3)^2)}+C
\end{equation}
$\Lambda$ has gauge freedom for constant values so we'll choose $C=0$ and get
\begin{equation}
  \Lambda = \frac{1}{2\alpha}e^{-\alpha((x^0)^2+(x^1)^2+(x^2)^2+(x^3)^2)}
\end{equation}
Next we'll find all the derivatives of $\Lambda$:
\begin{equation}
  \partial_\mu\Lambda = -x_\mu e^{-\alpha((x^0)^2+(x^1)^2+(x^2)^2+(x^3)^2)}
\end{equation}
raising the index
\begin{equation}
  \partial^\mu \Lambda = \mqty(-x^0\\x^1\\x^2\\x^3)e^{-\alpha((x^0)^2+(x^1)^2+(x^2)^2+(x^3)^2)}
\end{equation}
so the potential is
\begin{align}
  A'^\mu=A^\mu+\partial^\mu\Lambda&=x^\mu e^{-\alpha((x^0)^2+(x^1)^2+(x^2)^2+(x^3)^2)}\\
  &+\mqty(-x^0\\x^1\\x^2\\x^3)e^{-\alpha((x^0)^2+(x^1)^2+(x^2)^2+(x^3)^2)}\\
  &=2x^ie^{-\alpha((x^0)^2+(x^1)^2+(x^2)^2+(x^3)^2)}
\end{align}
\subsection{Switching potential}
noticing that the new potential can be written as
\begin{equation}
  A^\mu=x^\mu e^{-\alpha\tensor{\eta}{_\alpha_\beta}x^\alpha x^\beta}
\end{equation}
and that this looks similar to the form of the gauge transformation we used in the previous subsection. We'll try to fit a gauge transformation of a similar form and see what we get. We'll try
\begin{equation}
  \Lambda = -\frac{1}{2\alpha}e^{-\alpha\tensor{\eta}{_\alpha_\beta}x^\alpha x^\beta}
\end{equation}
This gives us the following derivatives:
\begin{align}
	\partial_\mu\Lambda&= -\frac{1}{2\alpha}\tensor{\eta}{_\mu_\gamma}\partial^\gamma\Lambda\\
	&= \tensor{\eta}{_\mu_\gamma}x^\gamma e^{-\alpha\tensor{\eta}{_\alpha_\beta}x^\alpha x^\beta}\\
	&= x_\gamma e^{-\alpha\tensor{\eta}{_\alpha_\beta}x^\alpha x^\beta}=A_\gamma
\end{align}
We found that we can define a gauge transformation that will just cancel out the potential giving us
\begin{equation}
  A'_\mu=A_\mu+\partial_\mu\Lambda=0
\end{equation}
And since the potential is zero, the fields tensor will also be 0, and since it's invariant for gauge transformations, we find that
\begin{align}
	\bold E=0 && \bold B=0
\end{align}